\documentclass[11pt,a4paper]{article}

\usepackage[utf8]{inputenc}
\usepackage[T1]{fontenc}
\usepackage[italian]{babel}
\usepackage{amsmath,amssymb}
\usepackage{booktabs}
\usepackage[margin=2.5cm]{geometry}

% Sezioni e sottosezioni sono numerate; i titoli di livello inferiore no
\setcounter{secnumdepth}{2}

\newcommand{\file}[1]{\texttt{\detokenize{#1}}}
\newcommand{\rif}{[L]}                      % il documento sulla linearizzazione
\newcommand{\Dc}{\mathcal{D}}
\newcommand{\Pc}{\mathcal{P}}
\newcommand{\uno}{\mathbf{1}}

\title{Costruzione del problema quadratico\\nell'MPC di bilanciamento attivo}
\author{Blocco \file{Problema_QP} del modello Simulink \file{Bilanciamento_MPC.slx}}
\date{}

\begin{document}
\maketitle

\noindent
A ogni passo di controllo l'MPC di bilanciamento deve risolvere un problema di
ottimo. Questo documento descrive come il blocco \file{Problema_QP} lo scrive
nella forma che il solutore \texttt{quadprog} sa risolvere: un problema
quadratico (QP). È la continuazione di \file{Linearizzazione_MPC.pdf}, citato
nel seguito come \rif: quel documento spiega come si ottiene la predizione
linearizzata delle celle, questo spiega che cosa se ne fa. Si può leggere da
solo: i risultati di \rif\ che servono sono richiamati nella
sezione~\ref{sec:incognite}, i concetti sui problemi quadratici nella
sezione~\ref{sec:richiami}.

\section{Il quadro generale}\label{sec:quadro}

Il controllo predittivo (MPC, \emph{model predictive control}) lavora così. A
ogni passo di controllo, all'istante $k$:
\begin{enumerate}
  \item misura lo stato delle celle;
  \item con un modello predice come evolveranno nei prossimi $p$ passi, in
        funzione delle correnti di bilanciamento $u^1, \dots, u^N$;
  \item sceglie le correnti che rendono migliore la predizione, secondo un
        costo e nel rispetto di alcuni vincoli;
  \item le applica per un solo passo, poi ricomincia da capo con le nuove
        misure.
\end{enumerate}
Il punto 3 è un problema di ottimo, risolto di nuovo a ogni passo: deve quindi
poter essere risolto in modo rapido e sicuro. Nel sottosistema
\file{Controllo} il lavoro è diviso tra quattro blocchi, in catena:

\begin{center}
\begin{tabular}{lll}
\toprule
Blocco & Che cosa calcola & Dove è descritto \\
\midrule
\file{Cella_nominale} & il riferimento $r_1, \dots, r_p$ & \rif \\
\file{Predizione_linearizzata} & $\bar Y$, $\bar S$, $D$, $E$ & \rif \\
\file{Problema_QP} & $H$, $f$, $A$, $b$, $A_{eq}$, $b_{eq}$, $z_0$ & questo documento \\
\file{Solutore_QP} & la soluzione $z$, che contiene le correnti & sezioni~\ref{sec:attivi} e~\ref{sec:partenza} \\
\bottomrule
\end{tabular}
\end{center}

\noindent
Il blocco \file{Problema_QP} non risolve nulla: prepara i dati. Riceve dei
numeri (la predizione, il riferimento, la mossa precedente) e restituisce
altri numeri, disposti in matrici: quelle che descrivono il problema da
risolvere. Capire il blocco significa capire perché ogni numero finisce in
quella posizione di quella matrice. Il documento procede così:
\begin{itemize}
  \item la sezione~\ref{sec:richiami} richiama che cos'è un problema
        quadratico e come lo si scrive;
  \item la sezione~\ref{sec:controllo} presenta il problema di controllo del
        paper e i tre ostacoli che impediscono di trattarlo subito come un QP;
  \item le sezioni dalla~\ref{sec:incognite} alla~\ref{sec:rimanenti}
        costruiscono il QP un pezzo alla volta: incognite, vincolo di tensione,
        costo, vincoli rimanenti;
  \item la sezione~\ref{sec:completo} mette insieme i pezzi in una procedura
        passo per passo, la \ref{sec:partenza} costruisce il punto di
        partenza, la \ref{sec:proprieta} ricava le proprietà del problema;
  \item le ultime tre sezioni trattano il caso in cui si bilancia il SOC,
        svolgono un esempio numerico e indicano dove si trova ogni passaggio
        nel codice.
\end{itemize}

\paragraph{Convenzioni.} In tutto il documento valgono le convenzioni
seguenti.
\begin{itemize}
  \item \emph{Numeri di equazione.} ``eq.~6'' e ``Remark~3'' rimandano al
        paper di Chen, Behal e Li (IEEE CDC 2021). I numeri tra sole
        parentesi, come (9), sono le equazioni di questo documento. Con
        \rif\ si rimanda a \file{Linearizzazione_MPC.pdf}: per esempio
        [L,~(6)] è la sua equazione (6).
  \item \emph{Indici.} $n = 1,\dots,N$ è la cella, scritta come indice in
        alto ($u^n$); $k$ è l'istante attuale; $j = 1,\dots,p$ è il passo lungo
        l'orizzonte. Le matrici $\bar Y$, $D$, $Y^0$ sono $p \times N$: la riga
        $j$ corrisponde all'istante $k+j$, la colonna $n$ alla cella $n$.
  \item \emph{Componenti.} $z_i$ è la componente $i$ del vettore $z$, $H_{il}$
        l'elemento di riga $i$ e colonna $l$ della matrice $H$; ${}^\top$
        indica la trasposizione.
  \item \emph{Barra.} Una grandezza con la barra, come $\bar u$ o $\bar Y$, è
        calcolata nel punto di linearizzazione, cioè con le correnti di
        bilanciamento del passo precedente.
  \item \emph{Correnti.} Sono positive in scarica. Una corrente di
        bilanciamento $u^n > 0$ significa quindi che il bilanciatore preleva
        carica dalla cella $n$; $u^n < 0$ che gliene cede.
\end{itemize}

\section{Richiami sui problemi quadratici}\label{sec:richiami}

\subsection{Definizione}

Un problema quadratico è un problema di minimo in cui il costo è un polinomio
di secondo grado nelle incognite e i vincoli sono di primo grado. Nella forma
usata da \texttt{quadprog}:
\begin{equation}
  \min_{z} \;\; \tfrac12\, z^\top H z + f^\top z
  \qquad \text{con} \qquad A\,z \le b , \quad A_{eq}\, z = b_{eq} .
  \label{eq:qp}
\end{equation}
Gli oggetti che vi compaiono sono:
\begin{itemize}
  \item $z \in \mathbb{R}^{n_z}$, il vettore colonna delle incognite
        $z_1, \dots, z_{n_z}$;
  \item $H$, matrice simmetrica $n_z \times n_z$, e $f \in \mathbb{R}^{n_z}$,
        che insieme definiscono il costo;
  \item $A$, matrice $m \times n_z$, e $b \in \mathbb{R}^{m}$, che definiscono
        $m$ vincoli di disuguaglianza;
  \item $A_{eq}$ e $b_{eq}$, che definiscono allo stesso modo i vincoli di
        uguaglianza.
\end{itemize}
Le disuguaglianze si leggono riga per riga: la riga $q$ dice che
\[
  A_{q1}\, z_1 + A_{q2}\, z_2 + \dots + A_{q n_z}\, z_{n_z} \;\le\; b_q .
\]
Ogni riga di $A$ è quindi un vincolo e ogni colonna un'incognita: l'elemento
$A_{qi}$ è il coefficiente con cui l'incognita $z_i$ compare nel vincolo $q$,
ed è zero se non vi compare.

\subsection{Come si leggono $H$ e $f$ da un costo}\label{sec:regola-costo}

Svolgendo i prodotti nella \eqref{eq:qp}, e usando la simmetria di $H$, il
costo si scrive
\begin{equation}
  \tfrac12\, z^\top H z + f^\top z
  = \tfrac12 \sum_{i} H_{ii}\, z_i^2
    + \sum_{i < l} H_{il}\, z_i z_l
    + \sum_{i} f_i\, z_i .
  \label{eq:scalare}
\end{equation}
Confrontando con un costo scritto per esteso si ottiene la regola che sarà
usata in tutto il documento:
\begin{itemize}
  \item se nel costo compare $a\,z_i^2$, allora $H_{ii} = 2a$;
  \item se compare $c\,z_i z_l$ con $i \ne l$, allora $H_{il} = H_{li} = c$;
  \item se compare $g\,z_i$, allora $f_i = g$;
  \item i termini che non contengono incognite si scartano: sommare una
        costante al costo non sposta il punto in cui è minimo.
\end{itemize}
Per esempio, il costo $3 z_1^2 + z_1 z_2 + 5 z_2^2 - 4 z_1 + 7$ corrisponde a
\[
  H = \begin{bmatrix} 6 & 1 \\ 1 & 10 \end{bmatrix}, \qquad
  f = \begin{bmatrix} -4 \\ 0 \end{bmatrix} .
\]
Il fattore $\tfrac12$ è una convenzione: fa sì che il gradiente del costo sia
$Hz + f$ e che $H$ sia la matrice delle sue derivate seconde. In questo
documento i prodotti misti $z_i z_l$ non compariranno mai: $H$ sarà sempre una
matrice diagonale.

\subsection{Come si scrive un vincolo}\label{sec:regola-vincolo}

Per portare un vincolo di primo grado nella forma di una riga di $A z \le b$
bastano tre mosse:
\begin{enumerate}
  \item portare a sinistra i termini con le incognite e a destra i numeri;
  \item se il verso è $\ge$, moltiplicare i due membri per $-1$, il che lo
        rovescia in $\le$;
  \item spezzare un vincolo doppio in due righe: $l \le z_i \le h$ diventa
        $z_i \le h$ e $-z_i \le -l$.
\end{enumerate}
Per esempio $2 z_1 \ge 3 + z_2$ diventa prima $2 z_1 - z_2 \ge 3$ e poi
$-2 z_1 + z_2 \le -3$: la riga di $A$ è $[\,-2 \;\; 1\,]$ e l'elemento di $b$ è
$-3$. Un vincolo di uguaglianza diventa, allo stesso modo, una riga di
$A_{eq}\, z = b_{eq}$.

\subsection{Convessità: perché un QP si risolve bene}\label{sec:convessita}

Una matrice simmetrica $H$ è \emph{semidefinita positiva} se
$z^\top H z \ge 0$ per ogni $z$, e \emph{definita positiva} se
$z^\top H z > 0$ per ogni $z \ne 0$. Per una matrice diagonale il controllo è
immediato, perché $z^\top H z = \sum_i H_{ii}\, z_i^2$: è semidefinita positiva
se e solo se tutti gli elementi della diagonale sono maggiori o uguali a zero,
definita positiva se e solo se sono tutti maggiori di zero. Valgono allora i
fatti seguenti.
\begin{itemize}
  \item L'insieme dei punti che rispettano tutti i vincoli si chiama
        \emph{insieme ammissibile}. Ogni disuguaglianza di primo grado
        individua un semispazio e ogni uguaglianza un iperpiano: l'insieme
        ammissibile è la loro intersezione, un poliedro. È un insieme
        convesso: il segmento che unisce due punti ammissibili è fatto di
        punti ammissibili.
  \item Se $H$ è semidefinita positiva il costo è una funzione convessa, e il
        QP si dice convesso. In un QP convesso ogni minimo locale è anche
        globale: un solutore che si ferma in un punto da cui non riesce più a
        migliorare ha trovato la soluzione, non un minimo ``sbagliato''.
  \item Se $H$ è definita positiva e l'insieme ammissibile non è vuoto, la
        soluzione esiste ed è unica.
  \item Se $H$ è solo semidefinita positiva, la soluzione esiste purché
        l'insieme ammissibile non sia vuoto e il costo vi sia limitato
        inferiormente (teorema di Frank e Wolfe); in generale può non essere
        unica.
\end{itemize}
Sono queste proprietà a rendere il QP adatto a un controllo in tempo reale:
la soluzione trovata è quella giusta, e i tempi di calcolo sono contenuti e
prevedibili.

\subsection{Punti ammissibili, vincoli attivi e metodo active-set}\label{sec:attivi}

Un punto $z$ è \emph{ammissibile} se rispetta tutti i vincoli. Un vincolo di
disuguaglianza è \emph{attivo} in $z$ se vi vale con il segno di uguaglianza:
il punto sta sul bordo del semispazio corrispondente. Nella soluzione di un QP
alcuni vincoli sono attivi e gli altri no; se si sapesse in anticipo quali,
basterebbe risolvere un sistema di equazioni di primo grado.

Il metodo \emph{active-set}, che \texttt{quadprog} usa nel modello, lo scopre
per tentativi ordinati. Parte da un punto, tiene un elenco dei vincoli che
considera attivi, risolve il problema come se quelli fossero uguaglianze e a
ogni iterazione aggiunge o toglie un vincolo dall'elenco, finché non può più
migliorare. Ha bisogno di un punto da cui partire: per questo il blocco gliene
prepara uno, $z_0$, già ammissibile (sezione~\ref{sec:partenza}).

\section{Il problema di controllo}\label{sec:controllo}

\subsection{Incognite e vincoli}

All'istante $k$ le incognite sono le correnti di bilanciamento, una per cella:
$u = (u^1, \dots, u^N)$. Ciascuna è tenuta costante su tutto l'orizzonte di
predizione, quindi le incognite sono $N$ e non $pN$. Devono rispettare tre
vincoli (eq.~6d--6f):
\begin{align}
  & u_{\min} \le u^n \le u_{\max} ,
     && n = 1,\dots,N , \label{eq:limiti}\\
  & y^n_{k+j} \ge y_{\min} ,
     && j = 1,\dots,p , \quad n = 1,\dots,N , \label{eq:tensione}\\
  & \sum_{n=1}^{N} u^n = 0 . \label{eq:somma}
\end{align}
\begin{itemize}
  \item La \eqref{eq:limiti} esprime i limiti del convertitore di
        bilanciamento (nel progetto $\pm 2$~A).
  \item Nella \eqref{eq:tensione} $y^n_{k+j}$ è la tensione predetta della
        cella $n$ all'istante $k+j$. Sotto $y_{\min}$ (nel progetto 3,2~V) la
        cella cede, e con lei l'intero pacco, perché le celle sono in serie. È
        il vincolo che dà senso a tutto il resto: lo scopo del bilanciamento è
        tenere tutte le celle lontane da $y_{\min}$ il più a lungo possibile.
  \item La \eqref{eq:somma} dice che il bilanciatore sposta carica da una
        cella all'altra, senza crearne né consumarne: ciò che toglie a una
        cella lo cede a un'altra.
\end{itemize}
Le tensioni predette dipendono da $u$ attraverso il modello della cella
[L,~sezione~2]: per questo la \eqref{eq:tensione} è un vincolo sulle correnti.

\subsection{I tre costi}

Il paper propone tre modi di misurare quanto è buona una scelta di $u$. In
tutti compare una \emph{grandezza bilanciata} $\sigma$, che può essere la
tensione ($\sigma = y$) o lo stato di carica ($\sigma = s$);
$\sigma^n_{k+j}$ è il suo valore predetto per la cella $n$ all'istante $k+j$.
\begin{align}
  J_t(u) &= \sum_{j=1}^{p} \sum_{n=1}^{N} \bigl(\sigma^n_{k+j} - r_j\bigr)^2
            + \rho \sum_{n=1}^{N} (u^n)^2
     && \text{(eq.~7)} \label{eq:jt}\\
  J_m(u) &= -\sum_{j=1}^{p} \min_{n}\, \sigma^n_{k+j}
            + \rho \sum_{n=1}^{N} (u^n)^2
     && \text{(eq.~8)} \label{eq:jm}\\
  J_\Delta(u) &= \sum_{j=1}^{p} \Bigl(\max_{n}\, \sigma^n_{k+j}
                 - \min_{n}\, \sigma^n_{k+j}\Bigr)
            + \rho \sum_{n=1}^{N} (u^n)^2
     && \text{(eq.~11)} \label{eq:jd}
\end{align}
\begin{itemize}
  \item Con $J_t$ ogni cella insegue un riferimento comune $r_j$: la
        traiettoria di una \emph{cella nominale}, cioè di una cella ideale con
        i parametri senza deviazioni, che parte dalla media degli stati
        misurati e subisce la sola corrente di pacco (blocco
        \file{Cella_nominale}). Se tutte le celle seguono la stessa
        traiettoria, restano allineate. Nel paper il riferimento è indicato
        con $\sigma^0$; nell'eq.~7 è scritto $\sigma^0_{k+1}$, qui inteso
        $\sigma^0_{k+j}$.
  \item Con $J_m$ si alza il più possibile, a ogni passo, la cella più bassa:
        è quella che arriverà per prima a $y_{\min}$. Il segno meno trasforma
        ``rendere massimo'' in ``rendere minimo''.
  \item Con $J_\Delta$ si riduce il divario tra la cella più alta e la più
        bassa.
\end{itemize}
Il secondo addendo è lo stesso nelle tre formulazioni: penalizza le correnti
di bilanciamento grandi, che sprecano energia per effetto Joule. Nel paper è
$u^\top R\, u$ con $R$ una matrice di peso; nel progetto $R = \rho\, I$, con
$\rho > 0$ scelto per ciascuna formulazione e proporzionale a $p$, perché il
primo addendo è una somma di $p$ termini e cresce con l'orizzonte.

\subsection{Tre ostacoli, tre rimedi}\label{sec:ostacoli}

Scritto così, il problema non è un QP, per tre ragioni. Ognuna ha un rimedio,
e i tre rimedi insieme portano alla forma \eqref{eq:qp}.
\begin{center}
\begin{tabular}{p{6.6cm}p{5.2cm}l}
\toprule
Ostacolo & Rimedio & Sezione \\
\midrule
Le tensioni predette non sono funzioni di primo grado di $u$: il modello della
cella non è lineare &
si sostituisce la predizione con la sua retta tangente &
\ref{sec:incognite} \\[4pt]
Il vincolo \eqref{eq:tensione} può essere impossibile da rispettare, e allora
il problema non ha soluzione &
si aggiunge un'incognita $e$ che lo allenta, a caro prezzo &
\ref{sec:vincolo-tensione} \\[4pt]
Il minimo e il massimo tra le celle, in $J_m$ e $J_\Delta$, non sono
polinomi &
si aggiungono le incognite $\varepsilon$, che ne prendono il posto &
\ref{sec:jm} e \ref{sec:jd} \\
\bottomrule
\end{tabular}
\end{center}

\section{Incognite e predizione linearizzata}\label{sec:incognite}

\subsection{La predizione linearizzata}

A ogni passo di controllo il blocco \file{Predizione_linearizzata} fornisce,
per ogni cella $n$ e per ogni passo $j$, due numeri [L,~(6)]:
\begin{itemize}
  \item $\bar Y_{j,n}$, la tensione che la cella avrebbe all'istante $k+j$ se
        la sua corrente di bilanciamento restasse quella del passo precedente,
        $\bar u^n = u^n_{k-1}$;
  \item $D_{j,n}$, il \emph{coefficiente di sensibilità}: di quanti volt cambia
        quella tensione per ogni ampere in più di corrente di bilanciamento.
\end{itemize}
Con essi la tensione predetta è approssimata dalla retta
$y^n_{k+j}(u^n) \simeq \bar Y_{j,n} + D_{j,n}\,(u^n - \bar u^n)$. Raccogliendo
i termini che non contengono l'incognita:
\begin{equation}
  y^n_{k+j}(u^n) = Y^0_{j,n} + D_{j,n}\,u^n ,
  \qquad
  Y^0_{j,n} = \bar Y_{j,n} - D_{j,n}\,\bar u^n .
  \label{eq:pred}
\end{equation}
$Y^0_{j,n}$ è la \emph{risposta libera}: la tensione che la retta predice
quando la cella non viene bilanciata ($u^n = 0$). Il prodotto $D_{j,n}\,u^n$ è
la \emph{risposta forzata}: l'effetto della corrente di bilanciamento. Nella
\eqref{eq:pred}, e da qui in avanti, si scrive $=$ al posto di $\simeq$: il QP
è costruito sul modello linearizzato come se fosse esatto. Quanto
l'approssimazione sia buona è discusso in [L,~sezione~8]. Tre proprietà della
\eqref{eq:pred} reggono tutta la costruzione.
\begin{itemize}
  \item \emph{È di primo grado in $u^n$.} Elevata al quadrato dà un polinomio
        di secondo grado; inserita in una disuguaglianza la lascia di primo
        grado. È proprio ciò che la forma \eqref{eq:qp} richiede.
  \item \emph{La tensione della cella $n$ dipende solo dalla sua corrente
        $u^n$}, non da quelle delle altre celle. Nel costo non nasceranno
        quindi prodotti misti $u^n u^m$.
  \item \emph{$D_{j,n}$ è negativo}: una corrente di bilanciamento positiva
        scarica di più la cella e ne abbassa la tensione. Nell'esempio di
        [L,~sezione~5.3] vale circa $-9$~mV/A dopo un passo e $-15$~mV/A dopo
        30.
\end{itemize}
Lo stesso vale per il SOC, con $\bar S$ ed $E$ al posto di $\bar Y$ e $D$. Nel
progetto, come nelle simulazioni del paper, la grandezza bilanciata è la
tensione: il documento è scritto per questo caso, e la
sezione~\ref{sec:soc} indica che cosa cambia se si bilancia il SOC.

\subsection{Le incognite del QP}

Oltre alle correnti, il QP ha alcune incognite in più, che servono a superare
il secondo e il terzo ostacolo della sezione~\ref{sec:ostacoli}. Sono raccolte
tutte in un solo vettore:
\begin{equation}
  z = \begin{bmatrix} u \\ \varepsilon \\ e \end{bmatrix} ,
  \label{eq:z}
\end{equation}
dove $u \in \mathbb{R}^N$ sono le correnti di bilanciamento, $\varepsilon$ le
variabili ausiliarie del costo (sezioni~\ref{sec:jm} e~\ref{sec:jd}) ed $e$ è
una sola variabile, che allenta il vincolo di tensione
(sezione~\ref{sec:vincolo-tensione}). Il numero delle $\varepsilon$ dipende
dalla formulazione:

\begin{center}
\begin{tabular}{lccc}
\toprule
Formulazione & Variabili $\varepsilon$ & Numero di $\varepsilon$ & Numero di incognite $n_z$ \\
\midrule
$J_t$ & nessuna & 0 & $N + 1$ \\
$J_m$ & $\varepsilon_1, \dots, \varepsilon_p$ & $p$ & $N + p + 1$ \\
$J_\Delta$ & $\varepsilon_1, \dots, \varepsilon_p$ e $\varepsilon_{p+1}, \dots, \varepsilon_{2p}$ & $2p$ & $N + 2p + 1$ \\
\bottomrule
\end{tabular}
\end{center}

\noindent
La posizione di ogni incognita dentro $z$ segue l'ordine della
\eqref{eq:z}:
\[
  u^n \to z_n , \qquad
  \varepsilon_j \to z_{N+j} , \qquad
  \varepsilon_{p+j} \to z_{N+p+j} , \qquad
  e \to z_{n_z} .
\]
In tutti i casi le correnti occupano le prime $N$ posizioni ed $e$ l'ultima.
Le variabili $\varepsilon$ ed $e$ non sono grandezze da applicare al pacco:
sono incognite di servizio, introdotte per poter scrivere il problema nella
forma \eqref{eq:qp}. Dalla soluzione $z$ si applicano solo le prime $N$
componenti.

\subsection{Come si costruiscono le matrici}\label{sec:metodo}

Il metodo è lo stesso per tutte le matrici. Si parte da $H$, $f$, $A$ e $b$
pieni di zeri, e si esamina un pezzo del problema alla volta: un termine del
costo oppure un gruppo di vincoli. Ogni pezzo viene scritto per esteso nelle
incognite, e le due regole della sezione~\ref{sec:richiami} dicono quali
elementi riempie e con quale valore:
\begin{itemize}
  \item un termine del costo riempie elementi di $H$ e di $f$
        (sezione~\ref{sec:regola-costo}), nelle posizioni delle incognite che
        contiene;
  \item un vincolo riempie una riga nuova di $A$ e un elemento di $b$
        (sezione~\ref{sec:regola-vincolo}): i coefficienti vanno nelle colonne
        delle incognite che contiene.
\end{itemize}
In entrambi i casi serve solo sapere la posizione di ogni incognita in $z$,
data qui sopra. Le sezioni dalla~\ref{sec:vincolo-tensione}
alla~\ref{sec:rimanenti} fanno questo lavoro un pezzo alla volta, spiegando
da dove viene ciascuno; la sezione~\ref{sec:completo} raccoglie poi tutti gli
elementi in una procedura unica, passo per passo.

\section{Il vincolo di tensione minima}\label{sec:vincolo-tensione}

\subsection{Dal vincolo rigido al vincolo allentato}

Con la \eqref{eq:pred} il vincolo \eqref{eq:tensione} diventa
$Y^0_{j,n} + D_{j,n}\,u^n \ge y_{\min}$ per ogni $j$ e $n$: sono $pN$
disuguaglianze di primo grado. Può però accadere che nessuna $u$ ammissibile le
soddisfi tutte. Per esempio, con $y_{\min} = 3{,}2$~V, $D_{j,n} = -15$~mV/A e
correnti limitate a $\pm 2$~A, la tensione predetta di una cella si può alzare
al più di 30~mV rispetto alla risposta libera, e solo dando alla cella $-2$~A:
se la risposta libera vale 3,15~V, il vincolo è violato qualunque sia $u$. Un
QP con insieme ammissibile vuoto non ha soluzione: il controllore resterebbe
senza risposta proprio quando una cella sta per cedere.

Il rimedio (eq.~17) è aggiungere un'incognita $e \ge 0$ che allenta il vincolo,
e farla pagare nel costo:
\begin{equation}
  Y^0_{j,n} + D_{j,n}\,u^n + e \;\ge\; y_{\min} , \qquad e \ge 0 ,
  \qquad \text{con il termine } W e^2 \text{ aggiunto al costo.}
  \label{eq:morbido}
\end{equation}
Un vincolo di questo tipo si dice \emph{morbido}, in contrapposizione a quello
\emph{rigido} \eqref{eq:tensione}. Conviene fissare alcuni fatti.
\begin{itemize}
  \item \emph{Il problema ha sempre punti ammissibili.} Presa una $u$
        qualunque che rispetti \eqref{eq:limiti} e \eqref{eq:somma}, basta
        scegliere $e$ abbastanza grande perché la \eqref{eq:morbido} valga per
        tutte le coppie $(j,n)$.
  \item \emph{$e$ ha un significato preciso.} Per $e \ge 0$ il costo cresce
        con $e$, quindi nella soluzione $e$ ha il più piccolo valore
        compatibile con i vincoli:
        \[
          e = \max\Bigl(0,\; \max_{j,n}\,\bigl(y_{\min} - y^n_{k+j}\bigr)\Bigr) .
        \]
        È di quanti volt, nel caso peggiore, la predizione scende sotto
        $y_{\min}$; vale zero se non ci scende mai. Per questo basta una sola
        $e$ per tutte le celle e tutti i passi.
  \item \emph{$W$ è molto più grande di $\rho$.} Violare il vincolo deve
        costare molto più che usare corrente: l'MPC lo fa solo se non può
        evitarlo. Nel progetto $W = 10^3$, mentre $\rho$ va da $10^{-3}$ a
        $10^{-1}$.
  \item \emph{Il vincolo morbido non equivale del tutto a quello rigido.}
        Finché nella soluzione tutte le tensioni predette stanno sopra
        $y_{\min}$, $e = 0$ e la variabile non ha alcun effetto. Quando invece
        una tensione predetta si fermerebbe esattamente a $y_{\min}$, all'MPC
        conviene una violazione piccolissima: vicino a $e = 0$ la penalità
        $W e^2$ ha derivata nulla, mentre il resto del costo diminuisce in
        proporzione alla violazione. La violazione che ne risulta è, con buona
        approssimazione, inversamente proporzionale a $W$: è un altro motivo
        per sceglierlo grande.
  \item Il paper introduce $e$ solo quando il problema rigido risulta senza
        soluzione; nel modello, per semplicità, $e$ è sempre presente.
\end{itemize}

\subsection{Le righe di $A$ e di $b$}

Applicando alla \eqref{eq:morbido} le mosse della
sezione~\ref{sec:regola-vincolo}, cioè incognite a sinistra, numeri a destra e
cambio di segno per avere il verso $\le$:
\begin{equation}
  -D_{j,n}\,u^n - e \;\le\; Y^0_{j,n} - y_{\min} ,
  \qquad j = 1,\dots,p , \quad n = 1,\dots,N .
  \label{eq:riga-tensione}
\end{equation}
C'è un vincolo per ogni coppia $(j, n)$, quindi $pN$ righe. In ciascuna i
coefficienti non nulli sono due: $-D_{j,n}$ nella colonna di $u^n$ e $-1$ nella
colonna di $e$. Il termine noto è il margine con cui la risposta libera sta
sopra $y_{\min}$. Il vincolo $e \ge 0$ è scritto più avanti, con gli altri
vincoli semplici (sezione~\ref{sec:rimanenti}).

\subsection{Il contributo al costo}

Il termine $W e^2$ contiene il quadrato di un'incognita, con coefficiente $W$.
Per la regola della sezione~\ref{sec:regola-costo}, indicando con $H_{ee}$
l'elemento della diagonale di $H$ nella posizione di $e$, cioè l'ultimo,
$H_{n_z n_z}$,
\begin{equation}
  H_{ee} = 2W .
  \label{eq:peso-e}
\end{equation}

\section{Il costo}\label{sec:costo}

Restano da scrivere gli altri elementi di $H$ e il vettore $f$. Il costo è la
somma di tre parti: il termine sulle correnti, comune alle tre formulazioni;
il termine $W e^2$, appena visto; la parte propria di ciascuna formulazione.
Le formulazioni $J_m$ e $J_\Delta$ portano con sé anche dei vincoli, che sono
descritti insieme al costo da cui nascono.

\subsection{Il termine sulle correnti}

Il termine $\rho \sum_n (u^n)^2$ contiene il quadrato di ogni corrente, con
coefficiente $\rho$. Poiché $u^n$ occupa la posizione $n$ di $z$,
\begin{equation}
  H_{nn} = 2\rho , \qquad n = 1,\dots,N .
  \label{eq:peso-u}
\end{equation}
Oltre al significato fisico, questo termine ha un ruolo matematico: rende il
costo strettamente convesso nelle correnti, e quindi unica la soluzione
(sezione~\ref{sec:proprieta}).

\subsection{$J_t$: inseguimento della cella nominale}\label{sec:jt}

Sostituendo la \eqref{eq:pred} nella \eqref{eq:jt}, il termine di inseguimento
è
\[
  \sum_{j=1}^{p}\sum_{n=1}^{N} \bigl(Y^0_{j,n} + D_{j,n}\,u^n - r_j\bigr)^2 .
\]
Ogni addendo è il quadrato di un binomio. Chiamando
$c_{j,n} = Y^0_{j,n} - r_j$ lo \emph{scarto} tra la risposta libera e il
riferimento, cioè quanto la cella starebbe lontana dal riferimento senza
bilanciamento,
\[
  \bigl(D_{j,n}\,u^n + c_{j,n}\bigr)^2
  = D_{j,n}^2\,(u^n)^2 + 2\,D_{j,n}\,c_{j,n}\,u^n + c_{j,n}^2 .
\]

\paragraph{Un esempio: due celle, primo passo.} Con $N = 2$ il primo passo
dell'orizzonte, $j = 1$, porta nel costo due addendi, uno per cella:
\begin{align*}
  & \bigl(D_{1,1}\,u^1 + c_{1,1}\bigr)^2 + \bigl(D_{1,2}\,u^2 + c_{1,2}\bigr)^2 \\[2pt]
  & \qquad = D_{1,1}^2\,(u^1)^2 + D_{1,2}^2\,(u^2)^2
      + 2\,D_{1,1}\,c_{1,1}\,u^1 + 2\,D_{1,2}\,c_{1,2}\,u^2
      + \bigl(c_{1,1}^2 + c_{1,2}^2\bigr) .
\end{align*}
Nel risultato ci sono i quadrati delle due correnti, le due correnti al primo
grado e una costante. Manca il prodotto $u^1 u^2$, perché ciascun binomio
contiene una sola corrente. Le correnti $u^1$ e $u^2$ occupano le posizioni 1 e
2 di $z$, quindi per la regola della sezione~\ref{sec:regola-costo} questo
passo contribuisce a $H$ e a $f$ con
\[
  \begin{bmatrix} H_{11} & H_{12} \\ H_{21} & H_{22} \end{bmatrix}
  \;\leftarrow\;
  \begin{bmatrix} 2\,D_{1,1}^2 & 0 \\ 0 & 2\,D_{1,2}^2 \end{bmatrix} ,
  \qquad
  \begin{bmatrix} f_1 \\ f_2 \end{bmatrix}
  \;\leftarrow\;
  \begin{bmatrix} 2\,D_{1,1}\,c_{1,1} \\ 2\,D_{1,2}\,c_{1,2} \end{bmatrix} ,
\]
mentre la costante $c_{1,1}^2 + c_{1,2}^2$ si scarta. Con i numeri
dell'esempio della sezione~\ref{sec:esempio}, cioè $D_{1,1} = -0{,}010$,
$c_{1,1} = 0{,}0250$ per la cella 1 e $D_{1,2} = -0{,}011$,
$c_{1,2} = -0{,}0255$ per la cella 2, i quattro contributi valgono
\begin{align*}
  2\,D_{1,1}^2 &= 2 \cdot 0{,}010^2 = 2{,}00 \cdot 10^{-4} ,
  & 2\,D_{1,1}\,c_{1,1} &= 2 \cdot (-0{,}010) \cdot 0{,}0250 = -5{,}00 \cdot 10^{-4} , \\
  2\,D_{1,2}^2 &= 2 \cdot 0{,}011^2 = 2{,}42 \cdot 10^{-4} ,
  & 2\,D_{1,2}\,c_{1,2} &= 2 \cdot (-0{,}011) \cdot (-0{,}0255) = 5{,}61 \cdot 10^{-4} .
\end{align*}
Il secondo passo, $j = 2$, porta altri due addendi dello stesso tipo, con
$D_{2,n}$ e $c_{2,n}$ al posto di $D_{1,n}$ e $c_{1,n}$: i loro contributi si
\emph{sommano} a questi, nelle stesse posizioni. Lo stesso accade per ogni
passo fino a $p$. È da qui che nascono le somme su $j$ della formula generale.

\paragraph{In generale.} Sommando su $j$ gli addendi che riguardano la stessa
cella $n$, il quadrato
$(u^n)^2$ ha coefficiente $\sum_j D_{j,n}^2$ e $u^n$ ha coefficiente
$2\sum_j D_{j,n}\,c_{j,n}$; l'ultimo termine non contiene incognite e si
scarta. Aggiungendo il termine sulle correnti \eqref{eq:peso-u}, la regola
della sezione~\ref{sec:regola-costo} dà
\begin{equation}
  H_{nn} = 2\sum_{j=1}^{p} D_{j,n}^2 + 2\rho , \qquad
  f_n = 2\sum_{j=1}^{p} D_{j,n}\,\bigl(Y^0_{j,n} - r_j\bigr) ,
  \qquad n = 1,\dots,N .
  \label{eq:jt-hf}
\end{equation}
Non ci sono prodotti misti, perché ogni binomio contiene una sola incognita:
il blocco di $H$ relativo alle correnti è diagonale. $J_t$ non richiede
variabili $\varepsilon$.

\paragraph{Lettura della formula.} Dimenticando per un momento i vincoli, il
costo è la somma di $N$ parabole, una per cella, e ciascuna ha il minimo in
\[
  u^n = -\frac{f_n}{H_{nn}}
      = \frac{\sum_j D_{j,n}\,\bigl(r_j - Y^0_{j,n}\bigr)}
             {\sum_j D_{j,n}^2 + \rho} .
\]
È la soluzione di un problema ai \emph{minimi quadrati}: si hanno $p$ scarti
da colmare, $r_j - Y^0_{j,n}$, e una sola incognita, $u^n$, che al passo $j$
sposta la tensione di $D_{j,n}\,u^n$. Non potendo annullare tutti gli scarti
con un solo numero, si sceglie quello che rende minima la somma dei quadrati
di ciò che ne resta. Il peso $\rho$ al denominatore frena la corrente: più è
grande, più
$u^n$ si avvicina a zero. Anche il segno torna: se la risposta libera di una
cella sta sopra il riferimento ($Y^0_{j,n} > r_j$), essendo $D_{j,n} < 0$
risulta $u^n > 0$, cioè la cella viene scaricata di più e la sua tensione
scende verso il riferimento.

I vincoli modificano poi questo valore ideale: la \eqref{eq:somma} lega tra
loro le correnti delle celle, che altrimenti sarebbero indipendenti, e la
\eqref{eq:limiti} le limita.

\subsection{$J_m$: massimizzazione della cella più bassa}\label{sec:jm}

Con la \eqref{eq:pred} il costo \eqref{eq:jm} contiene, per ogni passo $j$, il
termine $-\min_n \bigl(Y^0_{j,n} + D_{j,n}\,u^n\bigr)$. Il minimo tra $N$
funzioni di primo grado non è un polinomio: è una funzione lineare a tratti,
con uno spigolo ogni volta che cambia la cella più bassa. Non rientra quindi
nella forma \eqref{eq:qp}.

Il rimedio è introdurre per ogni passo $j$ una nuova incognita
$\varepsilon_j$, e chiedere che stia sotto le tensioni di tutte le celle:
\begin{equation}
  \min_{u,\,\varepsilon,\,e} \;\;
     -\sum_{j=1}^{p} \varepsilon_j + \rho \sum_{n=1}^{N} (u^n)^2 + W e^2
  \qquad \text{con} \qquad
  \varepsilon_j \le y^n_{k+j}
  \quad \text{per ogni } j \text{ e } n ,
  \label{eq:jm-eps}
\end{equation}
oltre ai vincoli già visti. $\varepsilon_j$ è un \emph{minorante} delle
tensioni al passo $j$: un livello sotto il quale non scende nessuna cella.

\paragraph{Perché le due scritture sono equivalenti.} Si fissi una $u$
qualunque. I vincoli della \eqref{eq:jm-eps} dicono che $\varepsilon_j$ non
può superare nessuna delle $N$ tensioni
$y^1_{k+j}, \dots, y^N_{k+j}$, cioè che
$\varepsilon_j \le \min_n y^n_{k+j}$. D'altra parte il costo diminuisce
quando $\varepsilon_j$ aumenta. Il valore migliore è quindi il più grande
consentito:
\[
  \varepsilon_j = \min_{n}\, y^n_{k+j} .
\]
Con questa scelta il costo della \eqref{eq:jm-eps} coincide con $J_m(u)$, a
parte il termine $W e^2$. Rendere minimo il costo rispetto a tutte le
incognite equivale quindi a rendere minimo $J_m$ rispetto a $u$: le due
scritture hanno la stessa soluzione. In breve, il costo spinge
$\varepsilon_j$ verso l'alto e i vincoli lo fermano all'altezza della cella
più bassa. Basta il vincolo da un solo lato (Remark~2 del paper): non serve
imporre che $\varepsilon_j$ sia \emph{uguale} al minimo, perché ci pensa il
costo.

\paragraph{Righe di $A$ e di $b$.} Sostituendo la \eqref{eq:pred} nel vincolo
$\varepsilon_j \le y^n_{k+j}$ e portando le incognite a sinistra:
\begin{equation}
  \varepsilon_j - D_{j,n}\,u^n \;\le\; Y^0_{j,n} ,
  \qquad j = 1,\dots,p , \quad n = 1,\dots,N .
  \label{eq:riga-minorante}
\end{equation}
Sono altre $pN$ righe, ciascuna con due coefficienti non nulli: $+1$ nella
colonna di $\varepsilon_j$ e $-D_{j,n}$ nella colonna di $u^n$.

\paragraph{Elementi di $H$ e di $f$.} Nel costo ogni $\varepsilon_j$ compare
al primo grado con coefficiente $-1$, e non compare al quadrato:
\begin{equation}
  f_{N+j} = -1 , \qquad H_{N+j,\,N+j} = 0 , \qquad j = 1,\dots,p .
  \label{eq:jm-hf}
\end{equation}
Per le correnti resta solo la \eqref{eq:peso-u}, e $f_n = 0$. Nelle
$\varepsilon$ il costo è di primo grado: $H$ ha degli zeri sulla diagonale, ed
è semidefinita positiva ma non definita positiva.

\subsection{$J_\Delta$: minima differenza tra le celle}\label{sec:jd}

Il costo \eqref{eq:jd} contiene, per ogni passo, sia il minimo sia il massimo
tra le celle. Si procede allo stesso modo, con due famiglie di incognite: le
$\varepsilon_j$ già viste, minoranti, e altre $p$ incognite
$\varepsilon_{p+j}$, che devono stare \emph{sopra} tutte le tensioni e ne sono
quindi dei \emph{maggioranti}:
\begin{equation}
\begin{aligned}
  & \min_{u,\,\varepsilon,\,e} \;\;
     \sum_{j=1}^{p} \varepsilon_{p+j} - \sum_{j=1}^{p} \varepsilon_j
     + \rho \sum_{n=1}^{N} (u^n)^2 + W e^2 \\[4pt]
  & \text{con} \qquad
  \varepsilon_j \;\le\; y^n_{k+j} \;\le\; \varepsilon_{p+j}
  \quad \text{per ogni } j \text{ e } n .
\end{aligned}
\label{eq:jd-eps}
\end{equation}
L'equivalenza si dimostra come prima. Fissata $u$, il costo spinge
$\varepsilon_j$ verso l'alto e $\varepsilon_{p+j}$ verso il basso; i vincoli li
fermano rispettivamente sulla cella più bassa e sulla più alta. Nella soluzione
\[
  \varepsilon_j = \min_{n}\, y^n_{k+j} , \qquad
  \varepsilon_{p+j} = \max_{n}\, y^n_{k+j} ,
\]
e la differenza $\varepsilon_{p+j} - \varepsilon_j$ è il divario tra le celle
al passo $j$.

I vincoli sui minoranti sono le righe \eqref{eq:riga-minorante}. Quelli sui
maggioranti, $y^n_{k+j} \le \varepsilon_{p+j}$, con la \eqref{eq:pred}
diventano
\begin{equation}
  D_{j,n}\,u^n - \varepsilon_{p+j} \;\le\; -Y^0_{j,n} ,
  \qquad j = 1,\dots,p , \quad n = 1,\dots,N ,
  \label{eq:riga-maggiorante}
\end{equation}
altre $pN$ righe con due coefficienti non nulli: $D_{j,n}$ nella colonna di
$u^n$ e $-1$ nella colonna di $\varepsilon_{p+j}$. Nel costo i minoranti hanno
coefficiente $-1$ e i maggioranti $+1$:
\begin{equation}
  f_{N+j} = -1 , \qquad f_{N+p+j} = +1 , \qquad j = 1,\dots,p ,
  \label{eq:jd-hf}
\end{equation}
con gli elementi corrispondenti della diagonale di $H$ uguali a zero.

\section{I vincoli rimanenti}\label{sec:rimanenti}

Restano i vincoli che non dipendono dalla predizione. I limiti di corrente
\eqref{eq:limiti} sono vincoli doppi, e ognuno si spezza in due righe:
\begin{equation}
  u^n \le u_{\max} , \qquad -u^n \le -u_{\min} ,
  \qquad n = 1,\dots,N .
  \label{eq:riga-limiti}
\end{equation}
Sono $2N$ righe con un solo coefficiente non nullo, $+1$ oppure $-1$ nella
colonna di $u^n$. Il segno della variabile $e$ si scrive
\begin{equation}
  -e \le 0 ,
  \label{eq:riga-e}
\end{equation}
una riga con $-1$ nella colonna di $e$. Infine la \eqref{eq:somma} è l'unico
vincolo di uguaglianza:
\begin{equation}
  A_{eq} = \bigl[\; \underbrace{1 \;\; \cdots \;\; 1}_{N}
              \;\; \underbrace{0 \;\; \cdots \;\; 0}_{n_z - N} \;\bigr] ,
  \qquad b_{eq} = 0 .
  \label{eq:uguaglianza}
\end{equation}
\texttt{quadprog} accetterebbe i limiti delle incognite anche come dati a
parte; nel modello sono scritti come righe di $A$, il che è equivalente.

\section{Il problema completo}\label{sec:completo}

Le sezioni dalla~\ref{sec:vincolo-tensione} alla~\ref{sec:rimanenti} hanno
ricavato i pezzi del problema uno alla volta, ciascuno con la sua
giustificazione. Qui sono rimessi in fila, senza più giustificazioni, come una
procedura da eseguire: quali elementi di $H$, $f$, $A$ e $b$ si scrivono, in
quale ordine e con quale valore. È la stessa procedura che segue il codice.

\subsection{Dimensioni e posizioni}\label{sec:dimensioni}

Prima di riempire le matrici bisogna sapere quanto sono grandi e dove sta ogni
incognita. Indicando con $t = 1, 2, 3$ la formulazione ($J_t$, $J_m$,
$J_\Delta$), il numero di incognite e il numero di vincoli di disuguaglianza
sono
\begin{equation}
  n_z = N + (t-1)\,p + 1 , \qquad
  n_{dis} = t\,pN + 2N + 1 .
  \label{eq:ndis}
\end{equation}
Il primo conta le $N$ correnti, le $(t-1)\,p$ variabili $\varepsilon$ e la
variabile $e$. Il secondo conta $t$ blocchi di $pN$ vincoli (quello di tensione
minima, più uno per i minoranti e uno per i maggioranti quando ci sono), i
$2N$ limiti di corrente e il segno di $e$. Con i valori del progetto, $N = 3$
e $p = 30$: 4 incognite e 97 vincoli per $J_t$, 34 e 187 per $J_m$, 64 e 277
per $J_\Delta$.

La posizione di un'incognita nel vettore $z$ decide tre cose insieme: il
numero della sua colonna in $A$ e in $A_{eq}$, il numero del suo elemento in
$f$, e il numero del suo elemento sulla diagonale di $H$.

\begin{center}
\begin{tabular}{llll}
\toprule
Incognita & & Posizione in $z$ & Presente in \\
\midrule
corrente $u^n$ & $n = 1,\dots,N$ & $n$ & $J_t$, $J_m$, $J_\Delta$ \\
minorante $\varepsilon_j$ & $j = 1,\dots,p$ & $N + j$ & $J_m$, $J_\Delta$ \\
maggiorante $\varepsilon_{p+j}$ & $j = 1,\dots,p$ & $N + p + j$ & $J_\Delta$ \\
variabile $e$ & & $n_z$, l'ultima & $J_t$, $J_m$, $J_\Delta$ \\
\bottomrule
\end{tabular}
\end{center}

\subsection{Costruzione di $H$ e di $f$}\label{sec:passi-costo}

Si parte da una matrice $H$ di $n_z \times n_z$ zeri e da un vettore $f$ di
$n_z$ zeri. Si scrivono poi, nell'ordine, solo gli elementi elencati qui sotto;
tutti gli altri restano nulli.
\begin{enumerate}
  \item \emph{Penalità su $e$} (tutte le formulazioni). Si scrive l'elemento
        della diagonale nella posizione di $e$, secondo la \eqref{eq:peso-e}:
        \mbox{$H_{n_z n_z} = 2W$}.
  \item \emph{Peso sulle correnti} (tutte le formulazioni). Per ogni cella
        $n = 1,\dots,N$ si scrive, secondo la \eqref{eq:peso-u},
        \mbox{$H_{nn} = 2\rho$}.
  \item \emph{Inseguimento} (solo $J_t$). Per ogni cella $n = 1,\dots,N$ si
        calcolano le due somme sui passi della \eqref{eq:jt-hf}; la prima si
        aggiunge all'elemento già scritto al passo 2, la seconda va in $f$:
        \[
          H_{nn} = 2\rho + 2\sum_{j=1}^{p} D_{j,n}^2 , \qquad
          f_n = 2\sum_{j=1}^{p} D_{j,n}\,\bigl(Y^0_{j,n} - r_j\bigr) .
        \]
  \item \emph{Minoranti} (solo $J_m$ e $J_\Delta$). Per ogni passo
        $j = 1,\dots,p$ si scrive, secondo la \eqref{eq:jm-hf},
        \mbox{$f_{N+j} = -1$}.
  \item \emph{Maggioranti} (solo $J_\Delta$). Per ogni passo $j = 1,\dots,p$
        si scrive, secondo la \eqref{eq:jd-hf}, \mbox{$f_{N+p+j} = +1$}.
\end{enumerate}
Nessun passo scrive fuori dalla diagonale di $H$: la matrice è diagonale in
tutte e tre le formulazioni. Scrivendo per esteso la diagonale di $H$ e il
vettore $f$, il risultato è
\begin{align*}
  J_t: \quad
    & \operatorname{diag}(H) = \bigl(\, H_{11},\, \dots,\, H_{NN},\; 2W \,\bigr) ,
    && f = \bigl(\, f_1,\, \dots,\, f_N,\; 0 \,\bigr) , \\[6pt]
  J_m: \quad
    & \operatorname{diag}(H) = \bigl(\, \underbrace{2\rho,\, \dots,\, 2\rho}_{N},\;
        \underbrace{0,\, \dots,\, 0}_{p},\; 2W \,\bigr) ,
    && f = \bigl(\, \underbrace{0,\, \dots,\, 0}_{N},\;
        \underbrace{-1,\, \dots,\, -1}_{p},\; 0 \,\bigr) , \\[6pt]
  J_\Delta: \quad
    & \operatorname{diag}(H) = \bigl(\, \underbrace{2\rho,\, \dots,\, 2\rho}_{N},\;
        \underbrace{0,\, \dots,\, 0}_{2p},\; 2W \,\bigr) ,
    && f = \bigl(\, \underbrace{0,\, \dots,\, 0}_{N},\;
        \underbrace{-1,\, \dots,\, -1}_{p},\;
        \underbrace{+1,\, \dots,\, +1}_{p},\; 0 \,\bigr) ,
\end{align*}
dove per $J_t$ gli elementi $H_{nn}$ e $f_n$ sono quelli del passo 3.

\subsection{Costruzione di $A$ e di $b$}\label{sec:passi-vincoli}

Si parte da una matrice $A$ di $n_{dis} \times n_z$ zeri e da un vettore $b$ di
$n_{dis}$ zeri. I vincoli si scrivono uno dopo l'altro, ciascuno in una riga
nuova, cominciando dalla prima. Scrivere un vincolo nella riga $q$ significa
due cose: mettere i suoi coefficienti negli elementi $A_{q,i}$, dove $i$ è la
posizione delle incognite che vi compaiono, e mettere il suo termine noto in
$b_q$. I vincoli si susseguono in sei blocchi, in quest'ordine.
\begin{enumerate}
  \item \emph{Tensione minima} (tutte le formulazioni),
        \eqref{eq:riga-tensione}. Per ogni cella $n = 1,\dots,N$ e, al suo
        interno, per ogni passo $j = 1,\dots,p$, una riga con $-D_{j,n}$ nella
        colonna di $u^n$ e $-1$ nella colonna di $e$; in $b$ va
        $Y^0_{j,n} - y_{\min}$.
  \item \emph{Minoranti} (solo $J_m$ e $J_\Delta$),
        \eqref{eq:riga-minorante}. Con lo stesso doppio ciclo, una riga con
        $-D_{j,n}$ nella colonna di $u^n$ e $+1$ nella colonna di
        $\varepsilon_j$; in $b$ va $Y^0_{j,n}$.
  \item \emph{Maggioranti} (solo $J_\Delta$), \eqref{eq:riga-maggiorante}.
        Ancora lo stesso doppio ciclo, una riga con $D_{j,n}$ nella colonna di
        $u^n$ e $-1$ nella colonna di $\varepsilon_{p+j}$; in $b$ va
        $-Y^0_{j,n}$.
  \item \emph{Limite superiore delle correnti} (tutte le formulazioni),
        \eqref{eq:riga-limiti}. Per ogni cella $n$, una riga con $+1$ nella
        colonna di $u^n$; in $b$ va $u_{\max}$.
  \item \emph{Limite inferiore delle correnti} (tutte le formulazioni),
        \eqref{eq:riga-limiti}. Per ogni cella $n$, una riga con $-1$ nella
        colonna di $u^n$; in $b$ va $-u_{\min}$.
  \item \emph{Segno di $e$} (tutte le formulazioni), \eqref{eq:riga-e}. Una
        sola riga, con $-1$ nella colonna di $e$; in $b$ va $0$.
\end{enumerate}
Nei primi tre blocchi i vincoli sono quindi ordinati per cella: prima i $p$
passi della cella 1, poi quelli della cella 2, e così via. La tabella riassume
la procedura in formule: per ogni blocco dà il numero $q$ della riga in cui va
il vincolo, i coefficienti da scrivere in quella riga e il termine noto.

\begin{center}
\setlength{\tabcolsep}{4pt}
\begin{tabular}{lllll}
\toprule
Blocco & Riga $q$ & Colonna di $u^n$ & Altra colonna & Termine noto \\
\midrule
1. Tensione minima & $(n-1)\,p + j$ & $A_{q,n} = -D_{j,n}$ & $A_{q,n_z} = -1$ & $b_q = Y^0_{j,n} - y_{\min}$ \\
2. Minoranti & $pN + (n-1)\,p + j$ & $A_{q,n} = -D_{j,n}$ & $A_{q,N+j} = +1$ & $b_q = Y^0_{j,n}$ \\
3. Maggioranti & $2pN + (n-1)\,p + j$ & $A_{q,n} = D_{j,n}$ & $A_{q,N+p+j} = -1$ & $b_q = -Y^0_{j,n}$ \\
4. Limite superiore & $t\,pN + n$ & $A_{q,n} = +1$ & --- & $b_q = u_{\max}$ \\
5. Limite inferiore & $t\,pN + N + n$ & $A_{q,n} = -1$ & --- & $b_q = -u_{\min}$ \\
6. Segno di $e$ & $t\,pN + 2N + 1$ & --- & $A_{q,n_z} = -1$ & $b_q = 0$ \\
\bottomrule
\end{tabular}
\end{center}

\noindent
Nella tabella $t$ è ancora la formulazione: i blocchi 4, 5 e 6 vengono subito
dopo l'ultimo blocco da $pN$ righe presente, cioè dopo $t\,pN$ righe. L'ultima
riga scritta ha numero $t\,pN + 2N + 1 = n_{dis}$, come previsto dalla
\eqref{eq:ndis}. Nel codice non serve calcolare il numero di riga con queste
formule: basta un contatore che parte da zero e avanza di uno a ogni vincolo
scritto.

Resta il vincolo di uguaglianza \eqref{eq:uguaglianza}, che è una riga sola:
$A_{eq}$ ha $1$ nelle colonne delle correnti, cioè nelle prime $N$, e zero
nelle altre; $b_{eq} = 0$.

\subsection{Il risultato a colpo d'occhio}\label{sec:blocchi}

La matrice $A$ costruita in questo modo ha una struttura a blocchi, che si
vede bene con tre simboli. Sia $y^0$ il vettore di lunghezza $pN$ ottenuto
mettendo in fila le colonne di $Y^0$, cioè nell'ordine in cui sono scritti i
vincoli; $\Dc$ la matrice $pN \times N$ che ha, nella colonna $n$, i
coefficienti $D_{1,n}, \dots, D_{p,n}$ in corrispondenza delle righe della
cella $n$, e zeri altrove; $\Pc$ la matrice $pN \times p$ formata da $N$ copie
della matrice identità $I_p$, una sopra l'altra:
\[
  y^0 = \begin{bmatrix} Y^0_{\cdot,1} \\ \vdots \\ Y^0_{\cdot,N} \end{bmatrix},
  \qquad
  \Dc = \begin{bmatrix}
          D_{\cdot,1} & & \\
          & \ddots & \\
          & & D_{\cdot,N}
        \end{bmatrix},
  \qquad
  \Pc = \begin{bmatrix} I_p \\ \vdots \\ I_p \end{bmatrix},
\]
dove $D_{\cdot,n}$ e $Y^0_{\cdot,n}$ sono le colonne $n$ di $D$ e di $Y^0$.
Con $\uno$ si indica un vettore di tutti uni. Per $J_\Delta$, la formulazione
più ricca, le incognite sono nell'ordine le correnti $u$, i minoranti
$\varepsilon_{1..p}$, i maggioranti $\varepsilon_{p+1..2p}$ ed $e$, e i vincoli
di disuguaglianza sono
\[
  \underbrace{
  \bordermatrix{
      & u & \varepsilon_{1..p} & \varepsilon_{p+1..2p} & e \cr
    \text{1. tensione minima} & -\Dc & 0 & 0 & -\uno \cr
    \text{2. minoranti}       & -\Dc & \Pc & 0 & 0 \cr
    \text{3. maggioranti}     & \Dc & 0 & -\Pc & 0 \cr
    \text{4. limite superiore} & I_N & 0 & 0 & 0 \cr
    \text{5. limite inferiore} & -I_N & 0 & 0 & 0 \cr
    \text{6. segno di } e     & 0 & 0 & 0 & -1 \cr
  }}_{A}
  \; z \;\le\;
  \underbrace{
  \begin{bmatrix}
    y^0 - y_{\min}\,\uno \\ y^0 \\ -y^0 \\ u_{\max}\,\uno \\ -u_{\min}\,\uno \\ 0
  \end{bmatrix}}_{b} .
\]
Ogni riga di questo schema è uno dei sei blocchi della procedura, e ogni
colonna un gruppo di incognite. Per $J_m$ si tolgono la riga dei maggioranti e
la colonna di $\varepsilon_{p+1..2p}$; per $J_t$ anche la riga dei minoranti e
la colonna di $\varepsilon_{1..p}$. Due cose si vedono a colpo d'occhio: $A$ è
quasi tutta fatta di zeri, con al più due coefficienti non nulli per riga, e
le correnti delle diverse celle non compaiono mai nello stesso vincolo di
disuguaglianza. Le lega soltanto il vincolo di uguaglianza.

\subsection{Una forma condensata}

Nel QP compaiono come incognite solo le correnti, oltre alle variabili di
servizio: gli stati delle celle lungo l'orizzonte non sono incognite, perché
la \eqref{eq:pred} li ha già espressi in funzione di $u$. È la cosiddetta
forma \emph{condensata}. Il paper conta invece $(2p+1)N$ incognite per $J_t$:
alle $N$ correnti si aggiungono i due stati di ogni cella a ogni passo, cioè
$2pN$ numeri, legati alle correnti da vincoli di uguaglianza. Le due forme
hanno la stessa soluzione; quella condensata è più piccola.

\section{Il punto di partenza}\label{sec:partenza}

Il blocco prepara anche un punto $z_0$ da cui il solutore parte. È la mossa
del passo precedente, completata con i valori delle variabili di servizio che
la rendono ammissibile:
\begin{equation}
  u = \bar u , \qquad
  \varepsilon_j = \min_{n}\, \bar Y_{j,n} , \qquad
  \varepsilon_{p+j} = \max_{n}\, \bar Y_{j,n} , \qquad
  e = \max\Bigl(0,\; y_{\min} - \min_{j,n}\, \bar Y_{j,n}\Bigr) .
  \label{eq:z0}
\end{equation}

\paragraph{Perché è ammissibile.} Tutto dipende da un'osservazione: nel punto
$u = \bar u$ la predizione linearizzata vale esattamente $\bar Y$. Infatti per
la \eqref{eq:pred}
$Y^0_{j,n} + D_{j,n}\,\bar u^n = \bar Y_{j,n}$. I vincoli si controllano
allora uno per uno.
\begin{itemize}
  \item \emph{Tensione minima} \eqref{eq:morbido}: serve
        $\bar Y_{j,n} + e \ge y_{\min}$ per ogni coppia. Vale perché
        $e \ge y_{\min} - \min_{j,n} \bar Y_{j,n} \ge y_{\min} - \bar Y_{j,n}$.
  \item \emph{Minoranti e maggioranti}: $\varepsilon_j$ è la più bassa delle
        tensioni $\bar Y_{j,1}, \dots, \bar Y_{j,N}$, quindi non ne supera
        nessuna; $\varepsilon_{p+j}$ è la più alta, quindi nessuna la supera.
  \item \emph{Segno di $e$}: $e$ è il massimo tra zero e un altro numero.
  \item \emph{Limiti e somma delle correnti}: $\bar u$ è la soluzione del QP
        del passo precedente, che rispettava \eqref{eq:limiti} e
        \eqref{eq:somma}. Al primo passo del modello $\bar u = 0$, che le
        rispetta purché $u_{\min} \le 0 \le u_{\max}$.
\end{itemize}

\paragraph{A che cosa serve.} Ha due usi. Il primo è dare al metodo
active-set un punto da cui partire (sezione~\ref{sec:attivi}): un punto già
ammissibile, e di solito vicino alla soluzione, perché le correnti di
bilanciamento cambiano poco da un passo all'altro. Il secondo è di sicurezza:
se \texttt{quadprog} non restituisce una soluzione, il blocco
\file{Solutore_QP} restituisce $z_0$, cioè mantiene la mossa precedente.
``Continuare come prima'' è sempre una scelta ammissibile.

\section{Proprietà del problema}\label{sec:proprieta}

Con i richiami della sezione~\ref{sec:convessita} si può ora dire che tipo di
problema è stato costruito.
\begin{itemize}
  \item \emph{È convesso.} $H$ è diagonale e i suoi elementi sono
        $2\rho$ (più una somma di quadrati per $J_t$), zero oppure $2W$:
        nessuno è negativo, quindi $H$ è semidefinita positiva. Ogni minimo
        locale è globale.
  \item \emph{Ha punti ammissibili.} Uno è $z_0$.
  \item \emph{Ha una soluzione.} Per $J_t$ tutti gli elementi della diagonale
        sono positivi: $H$ è definita positiva e la soluzione esiste. Per
        $J_m$ e $J_\Delta$ $H$ ha degli zeri sulla diagonale, ma il costo è
        limitato inferiormente. Le correnti sono limitate dalla
        \eqref{eq:limiti}, quindi lo sono le tensioni predette: i minoranti
        $\varepsilon_j$ non possono salire oltre un certo valore e i
        maggioranti $\varepsilon_{p+j}$ non possono scendere sotto un altro.
        Gli addendi $-\varepsilon_j$ e $+\varepsilon_{p+j}$ hanno quindi un
        limite inferiore, e gli altri addendi del costo non sono mai negativi.
  \item \emph{La soluzione è unica.} Per $J_t$ segue dal fatto che $H$ è
        definita positiva. Per $J_m$ e $J_\Delta$ serve un argomento in più.
        Se ci fossero due soluzioni diverse, per la convessità sarebbero
        soluzioni anche tutti i punti del segmento che le unisce, e lungo il
        segmento il costo sarebbe costante. Ma lungo un segmento il costo è un
        polinomio di secondo grado nel parametro che lo percorre, e il
        coefficiente del termine di secondo grado è
        $\rho \sum_n (\Delta u^n)^2 + W\,(\Delta e)^2$, dove $\Delta u$ e
        $\Delta e$ sono le differenze tra le due soluzioni. Il costo può essere
        costante solo se questo coefficiente è nullo, cioè se le due soluzioni
        hanno le stesse correnti e la stessa $e$. Fissata $u$, d'altra parte, i
        valori migliori delle $\varepsilon$ sono determinati: il minimo e il
        massimo tra le celle. Le due soluzioni coincidono.
\end{itemize}
C'è infine una ragione per cui nel modello si usa il metodo active-set. In
$J_m$ e $J_\Delta$ il costo è di primo grado nelle $\varepsilon$, e nella
soluzione ogni $\varepsilon$ sta esattamente su un vincolo, cioè sulla tensione
di una cella. Un metodo che lavora con i vincoli attivi individua questi
vincoli in modo esatto, invece di avvicinarvisi soltanto.

\section{Se si bilancia il SOC}\label{sec:soc}

Tutto il documento è scritto per $\sigma = y$. Se la grandezza bilanciata è il
SOC, serve anche la sua predizione linearizzata, che arriva dallo stesso
blocco di quella della tensione:
\[
  s^n_{k+j}(u^n) = S^0_{j,n} + E_{j,n}\,u^n ,
  \qquad
  S^0_{j,n} = \bar S_{j,n} - E_{j,n}\,\bar u^n .
\]
Il riferimento $r_j$ diventa il SOC della cella nominale. Le sostituzioni da
fare sono due, e riguardano solo le parti che contengono $\sigma$.
\begin{itemize}
  \item \emph{Nel costo e nei vincoli sulle $\varepsilon$}, cioè nelle
        \eqref{eq:jt-hf}, \eqref{eq:riga-minorante} e
        \eqref{eq:riga-maggiorante}, al posto di $Y^0$ e $D$ si usano $S^0$ ed
        $E$. Nel punto di partenza \eqref{eq:z0} le $\varepsilon$ si calcolano
        da $\bar S$ invece che da $\bar Y$.
  \item \emph{Il vincolo di tensione minima} \eqref{eq:riga-tensione} e il
        valore di $e$ nella \eqref{eq:z0} restano com'erano: il vincolo
        \eqref{eq:tensione} riguarda sempre la tensione, qualunque sia la
        grandezza bilanciata.
\end{itemize}
In questo caso il blocco usa quindi tutte e quattro le matrici $\bar Y$,
$\bar S$, $D$, $E$. Va notato che per il SOC la retta non è
un'approssimazione: il SOC è esattamente di primo grado nella corrente, finché
la cella non è piena [L,~sezione~5].

\section{Esempio numerico}\label{sec:esempio}

Per vedere i simboli diventare numeri si considera un caso piccolo, con valori
scelti per fare i conti a mano: $N = 2$ celle e $p = 2$ passi. I dati sono
\[
  \bar Y = \begin{bmatrix} 3{,}700 & 3{,}660 \\ 3{,}690 & 3{,}648 \end{bmatrix}
  \text{V} , \quad
  D = \begin{bmatrix} -0{,}010 & -0{,}011 \\ -0{,}012 & -0{,}013 \end{bmatrix}
  \frac{\text{V}}{\text{A}} , \quad
  \bar u = \begin{bmatrix} 0{,}5 \\ -0{,}5 \end{bmatrix} \text{A} , \quad
  r = \begin{bmatrix} 3{,}680 \\ 3{,}669 \end{bmatrix} \text{V} ,
\]
con $y_{\min} = 3{,}2$~V, $u_{\min} = -2$~A, $u_{\max} = 2$~A e $W = 1000$. La
cella 1 (prima colonna) è la più alta, la cella 2 la più bassa.

\paragraph{Risposta libera.} Dalla \eqref{eq:pred},
$Y^0_{j,n} = \bar Y_{j,n} - D_{j,n}\,\bar u^n$:
\[
  Y^0 = \begin{bmatrix}
          3{,}700 + 0{,}010 \cdot 0{,}5 & 3{,}660 - 0{,}011 \cdot 0{,}5 \\
          3{,}690 + 0{,}012 \cdot 0{,}5 & 3{,}648 - 0{,}013 \cdot 0{,}5
        \end{bmatrix}
      = \begin{bmatrix} 3{,}7050 & 3{,}6545 \\ 3{,}6960 & 3{,}6415 \end{bmatrix}
      \text{V} .
\]
Senza bilanciamento la cella 1 starebbe un po' più in alto di quanto dice
$\bar Y$, e la cella 2 un po' più in basso: la mossa precedente stava già
avvicinando le due celle.

\paragraph{$J_t$: costruzione di $H$ e di $f$.} Le incognite sono
$z = (u^1, u^2, e)$: $n_z = 3$, le correnti occupano le posizioni 1 e 2 ed
$e$ la posizione 3. Si seguono i passi della sezione~\ref{sec:passi-costo},
con $\rho = 10^{-4}$.
\begin{enumerate}
  \item Penalità su $e$: $H_{33} = 2W = 2000$.
  \item Peso sulle correnti: $H_{11} = H_{22} = 2\rho = 2 \cdot 10^{-4}$.
  \item Inseguimento: a ciascuno dei due elementi si aggiunge la somma sui
        passi, e si calcolano $f_1$ e $f_2$:
        \begin{align*}
          H_{11} &= 2 \cdot 10^{-4} + 2\,(0{,}010^2 + 0{,}012^2)
                  = 6{,}88 \cdot 10^{-4} , \\
          H_{22} &= 2 \cdot 10^{-4} + 2\,(0{,}011^2 + 0{,}013^2)
                  = 7{,}80 \cdot 10^{-4} , \\
          f_1 &= 2\,\bigl[(-0{,}010) \cdot 0{,}0250 + (-0{,}012) \cdot 0{,}0270\bigr]
                  = -1{,}148 \cdot 10^{-3} , \\
          f_2 &= 2\,\bigl[(-0{,}011) \cdot (-0{,}0255) + (-0{,}013) \cdot (-0{,}0275)\bigr]
                  = 1{,}276 \cdot 10^{-3} .
        \end{align*}
\end{enumerate}
I numeri $0{,}0250$ e $0{,}0270$ sono gli scarti $Y^0_{j,1} - r_j$ della cella
1, che sta sopra il riferimento; quelli della cella 2, che sta sotto, sono
$-0{,}0255$ e $-0{,}0275$. I passi 4 e 5 non riguardano $J_t$. In tutto:
\[
  H = \begin{bmatrix}
        6{,}88 \cdot 10^{-4} & 0 & 0 \\
        0 & 7{,}80 \cdot 10^{-4} & 0 \\
        0 & 0 & 2000
      \end{bmatrix} ,
  \qquad
  f = \begin{bmatrix} -1{,}148 \cdot 10^{-3} \\ 1{,}276 \cdot 10^{-3} \\ 0 \end{bmatrix} .
\]

\paragraph{$J_t$: costruzione di $A$ e di $b$.} I vincoli di disuguaglianza
sono $n_{dis} = 1 \cdot 4 + 4 + 1 = 9$. Seguendo i blocchi della
sezione~\ref{sec:passi-vincoli}: il blocco 1 occupa le righe da 1 a 4, i
blocchi 4 e 5 le righe da 5 a 8, il blocco 6 la riga 9. I blocchi 2 e 3 non ci
sono.
\[
  \bordermatrix{
      & u^1 & u^2 & e \cr
    \text{1:\; tensione, cella 1, } j = 1 & 0{,}010 & 0 & -1 \cr
    \text{2:\; tensione, cella 1, } j = 2 & 0{,}012 & 0 & -1 \cr
    \text{3:\; tensione, cella 2, } j = 1 & 0 & 0{,}011 & -1 \cr
    \text{4:\; tensione, cella 2, } j = 2 & 0 & 0{,}013 & -1 \cr
    \text{5:\; limite superiore, cella 1} & 1 & 0 & 0 \cr
    \text{6:\; limite superiore, cella 2} & 0 & 1 & 0 \cr
    \text{7:\; limite inferiore, cella 1} & -1 & 0 & 0 \cr
    \text{8:\; limite inferiore, cella 2} & 0 & -1 & 0 \cr
    \text{9:\; segno di } e & 0 & 0 & -1 \cr
  }
  \begin{bmatrix} u^1 \\ u^2 \\ e \end{bmatrix}
  \;\le\;
  \begin{bmatrix}
    0{,}5050 \\ 0{,}4960 \\ 0{,}4545 \\ 0{,}4415 \\ 2 \\ 2 \\ 2 \\ 2 \\ 0
  \end{bmatrix} .
\]
Il numero davanti a ogni riga è il numero $q$ della
sezione~\ref{sec:passi-vincoli}: per esempio il vincolo di tensione della
cella 2 al passo 1 sta nella riga $(n-1)\,p + j = (2-1) \cdot 2 + 1 = 3$. Nelle
prime quattro righe i coefficienti $0{,}010$, $0{,}012$, \dots\ sono
$-D_{j,n}$, scritti nella colonna della cella a cui si riferiscono, e i termini
noti sono i margini $Y^0_{j,n} - 3{,}2$. Il vincolo di uguaglianza è
$A_{eq} = [\,1 \;\; 1 \;\; 0\,]$, $b_{eq} = 0$, e il punto di partenza è
$z_0 = (0{,}5;\; -0{,}5;\; 0)$.

\paragraph{La soluzione per $J_t$.} Senza vincoli le due parabole avrebbero il
minimo in $u^1 = -f_1/H_{11} = 1{,}669$~A e $u^2 = -f_2/H_{22} = -1{,}636$~A:
la cella alta andrebbe scaricata di più, la bassa di meno. La somma però non
è nulla. Imponendo $u^2 = -u^1$ e rendendo minimo il costo rispetto alla sola
$u^1$ si ottiene
\[
  u^1 = \frac{f_2 - f_1}{H_{11} + H_{22}} = 1{,}651~\text{A} , \qquad
  u^2 = -1{,}651~\text{A} , \qquad e = 0 ,
\]
che rispetta i limiti di $\pm 2$~A. Le tensioni predette con queste correnti
vanno da 3,663 a 3,688~V: sono lontane da $y_{\min}$, quindi il vincolo di
tensione non interviene ed $e = 0$. È il risultato che restituisce
\texttt{quadprog}.

\paragraph{Che cosa cambia per $J_m$.} Le incognite diventano
$z = (u^1, u^2, \varepsilon_1, \varepsilon_2, e)$: $n_z = 5$, i minoranti
occupano le posizioni 3 e 4 ed $e$ la posizione 5. Per $H$ e $f$ si eseguono
i passi 1, 2 e 4 della sezione~\ref{sec:passi-costo}:
\[
  \operatorname{diag}(H) = \bigl(\, 2\rho,\; 2\rho,\; 0,\; 0,\; 2000 \,\bigr) ,
  \qquad
  f = \bigl(\, 0,\; 0,\; -1,\; -1,\; 0 \,\bigr) .
\]
Per $A$ e $b$ compare il blocco 2, subito dopo quello di tensione. Le prime
quattro righe sono quelle di prima, con il $-1$ di $e$ che ora sta nella
quinta colonna; seguono le righe da 5 a 8,
\[
  \bordermatrix{
      & u^1 & u^2 & \varepsilon_1 & \varepsilon_2 & e \cr
    \text{5:\; minorante, cella 1, } j = 1 & 0{,}010 & 0 & 1 & 0 & 0 \cr
    \text{6:\; minorante, cella 1, } j = 2 & 0{,}012 & 0 & 0 & 1 & 0 \cr
    \text{7:\; minorante, cella 2, } j = 1 & 0 & 0{,}011 & 1 & 0 & 0 \cr
    \text{8:\; minorante, cella 2, } j = 2 & 0 & 0{,}013 & 0 & 1 & 0 \cr
  }
  \; z \;\le\;
  \begin{bmatrix} 3{,}7050 \\ 3{,}6960 \\ 3{,}6545 \\ 3{,}6415 \end{bmatrix} ,
\]
e poi i limiti di corrente, che scendono alle righe da 9 a 12, e il segno di
$e$, alla riga 13. In ogni riga del blocco 2 il coefficiente $+1$ sta nella
colonna del minorante del passo corrispondente: $\varepsilon_1$ per $j = 1$,
$\varepsilon_2$ per $j = 2$. Il punto di partenza è
$z_0 = (0{,}5;\; -0{,}5;\; 3{,}660;\; 3{,}648;\; 0)$: le due $\varepsilon$
sono le tensioni $\bar Y$ della cella più bassa.

\paragraph{Che cosa cambia per $J_\Delta$.} Le incognite diventano
$z = (u^1, u^2, \varepsilon_1, \varepsilon_2, \varepsilon_3, \varepsilon_4, e)$:
$n_z = 7$, i maggioranti occupano le posizioni 5 e 6 ed $e$ la posizione 7.
Ai passi di $J_m$ si aggiunge il passo 5, quindi
$\operatorname{diag}(H) = (\, 2\rho,\, 2\rho,\, 0,\, 0,\, 0,\, 0,\, 2000 \,)$ e
$f = (\, 0,\, 0,\, -1,\, -1,\, +1,\, +1,\, 0 \,)$. In $A$ compare anche il
blocco 3, nelle righe da 9 a 12,
\[
  \bordermatrix{
      & u^1 & u^2 & \varepsilon_1 & \varepsilon_2 & \varepsilon_3 & \varepsilon_4 & e \cr
    \text{9:\; maggiorante, cella 1, } j = 1 & -0{,}010 & 0 & 0 & 0 & -1 & 0 & 0 \cr
    \text{10:\; maggiorante, cella 1, } j = 2 & -0{,}012 & 0 & 0 & 0 & 0 & -1 & 0 \cr
    \text{11:\; maggiorante, cella 2, } j = 1 & 0 & -0{,}011 & 0 & 0 & -1 & 0 & 0 \cr
    \text{12:\; maggiorante, cella 2, } j = 2 & 0 & -0{,}013 & 0 & 0 & 0 & -1 & 0 \cr
  }
  \; z \;\le\;
  \begin{bmatrix} -3{,}7050 \\ -3{,}6960 \\ -3{,}6545 \\ -3{,}6415 \end{bmatrix} ,
\]
mentre i limiti di corrente scendono alle righe da 13 a 16 e il segno di $e$
alla riga 17. Il punto di partenza è
$z_0 = (0{,}5;\; -0{,}5;\; 3{,}660;\; 3{,}648;\; 3{,}700;\; 3{,}690;\; 0)$.

\paragraph{Le tre soluzioni a confronto.} Come nel progetto, il peso $\rho$ è
diverso per le tre formulazioni: $10^{-4}$ per $J_t$, $5 \cdot 10^{-3}$ per
$J_m$, $10^{-2}$ per $J_\Delta$. Le soluzioni di \texttt{quadprog} sono:

\begin{center}
\begin{tabular}{lccccc}
\toprule
 & $u^1$ [A] & $u^2$ [A] & $\varepsilon_1$;\; $\varepsilon_2$ [V] & $\varepsilon_3$;\; $\varepsilon_4$ [V] & $e$ [V] \\
\midrule
$J_t$ & $\phantom{-}1{,}651$ & $-1{,}651$ & --- & --- & 0 \\
$J_m$ & $\phantom{-}1{,}200$ & $-1{,}200$ & 3,66770;\; 3,65710 & --- & 0 \\
$J_\Delta$ & $\phantom{-}1{,}150$ & $-1{,}150$ & 3,66715;\; 3,65645 & 3,69350;\; 3,68220 & 0 \\
\bottomrule
\end{tabular}
\end{center}

\noindent
In tutti e tre i casi $e = 0$: le tensioni predette sono lontane da
$y_{\min}$ e il vincolo di tensione non interviene. Il bilanciatore preleva
corrente dalla cella alta e la cede a quella bassa, in quantità diverse perché
diversi sono i criteri.

Nelle ultime due righe si riconosce il meccanismo delle sezioni~\ref{sec:jm}
e~\ref{sec:jd}: i minoranti si sono fermati sulle tensioni della cella più
bassa, $\varepsilon_j = Y^0_{j,2} + D_{j,2}\,u^2$, e i maggioranti su quelle
della più alta. Per $J_m$, per esempio,
$\varepsilon_1 = 3{,}6545 + 0{,}011 \cdot 1{,}2 = 3{,}6677$~V. Nel linguaggio
della sezione~\ref{sec:attivi}, i vincoli sui minoranti della cella 2 sono
attivi, mentre quelli della cella 1, che sta più in alto, non lo sono; per
$J_\Delta$ sono attivi anche i vincoli sui maggioranti della cella 1. In tutti
e tre i casi è attivo il vincolo $e \ge 0$, e nessun limite di corrente.

\section{Dove si trova nel codice}\label{sec:codice}

Il blocco \file{Problema_QP} chiama una funzione che riceve la predizione e
restituisce le matrici. Ne esistono due versioni, con gli stessi ingressi e le
stesse uscite:
\begin{itemize}
  \item \file{qp_problema.m} riempie le matrici a blocchi, con operazioni su
        vettori;
  \item \file{problema_qp_semplice.m} le riempie un elemento alla volta, per
        una cella e un passo per volta, ed è diviso nelle stesse parti di
        questo documento: è la versione da leggere.
\end{itemize}
Le due funzioni producono le stesse matrici, e nel blocco si possono
scambiare: per usare la versione da leggere basta che il blocco contenga
l'istruzione
{\small
\begin{verbatim}
    [H, f, A, b, Aeq, beq, z0] = ...
        problema_qp_semplice(sl, sigma0, Y, S, dY, dS, u_prec);
\end{verbatim}
}
\noindent
dove \texttt{sigma0}, \texttt{dY}, \texttt{dS} e \file{u_prec} sono i nomi che
il riferimento $r$, le matrici $D$ ed $E$ e la mossa $\bar u$ hanno nel
modello. La versione da leggere esegue la procedura della
sezione~\ref{sec:completo} nello stesso ordine: ogni vincolo è scritto in una
riga nuova di \texttt{A}, e la variabile \texttt{q} è il contatore di righe
della sezione~\ref{sec:passi-vincoli}, che avanza di uno a ogni vincolo. Le
variabili hanno il nome dei simboli del documento:

\begin{center}
\begin{tabular}{ll@{\qquad}ll}
\toprule
Simbolo & Nome nel codice & Simbolo & Nome nel codice \\
\midrule
$\bar Y$, $D$ & \texttt{Y}, \texttt{D} & $H$, $f$ & \texttt{H}, \texttt{f} \\
$\bar S$, $E$ & \texttt{S}, \texttt{E} & $A$, $b$ & \texttt{A}, \texttt{b} \\
$Y^0$ & \texttt{Y0} & $A_{eq}$, $b_{eq}$ & \texttt{Aeq}, \texttt{beq} \\
$\bar u$ & \file{u_bar} & $z_0$ & \texttt{z0} \\
$r_j$ & \texttt{r(j)} & $n_z$, $n_{dis}$ & \file{n_z}, \file{n_dis} \\
$\rho$, $W$ & \texttt{rho}, \texttt{W} & posizione di $e$ in $z$ & \file{i_e} \\
$y_{\min}$, $u_{\min}$, $u_{\max}$ & \file{y_min}, \file{u_min}, \file{u_max} & formulazione $t$ & \texttt{tipo} \\
\bottomrule
\end{tabular}
\end{center}

\noindent
Nel costo e nei vincoli sulle $\varepsilon$ il codice usa \texttt{Yb},
\texttt{Db} e \texttt{Yb0} (``b'' per ``bilanciata''): coincidono con
\texttt{Y}, \texttt{D} e \texttt{Y0} se si bilancia la tensione, e con le
grandezze corrispondenti del SOC altrimenti (sezione~\ref{sec:soc}). I
parametri arrivano dalla struttura \texttt{sl}, creata da
\file{init_simulink.m}: $\rho$ è \texttt{sl.r}, la formulazione è
\texttt{sl.tipo}.

\end{document}
